\chapter{Conclusion}

This project demonstrates that a practical, local-first RAG system for campus information retrieval can be built entirely from open-weight models and open-source infrastructure, without relying on external API services. The system covers the complete pipeline from web crawling through text preprocessing, semantic chunking, hybrid retrieval, cross-encoder reranking, self-RAG reflection, and streaming generation---deployed on hardware ranging from a single CPU laptop to a three-GPU server.

Several design decisions are worth highlighting. The paragraph-first chunking strategy with structured-page detection preserves the coherence of short announcement-style content that would be mangled by naive fixed-window splitting. The hybrid retrieval architecture addresses a known weakness of dense-only retrieval on exact-match queries, which are common in the campus domain (specific course codes, building names, contact details). The self-RAG loop introduces explicit scope and factual-support checks that prevent the model from generating confident but ungrounded responses when the indexed corpus does not contain the answer. The relay-worker deployment decouples API availability from GPU hardware placement, enabling flexible distributed operation without requiring the GPU machine to be publicly addressable.

The most important gap identified by this work is the absence of quantitative evaluation. The design choices enumerated above are motivated by prior literature and engineering reasoning, but their actual effect on this system's answer quality over this corpus has not been measured. As detailed in Chapter~\ref{ch:limitations}, constructing a labelled benchmark for the SUSTech domain and using it to run ablations across retrieval modes, chunk sizes, and self-RAG configurations is the highest-priority next step. Without such evaluation, the system is a well-engineered prototype rather than a validated tool.

Subject to that caveat, the system is functional and deployable. It indexes SUSTech public web content, answers in-scope queries with source citations, explicitly acknowledges out-of-scope or insufficiently evidenced questions, and streams responses in real time to a responsive web interface.
